% !TeX root = ../main.tex
\chapter{Introduction}\label{ch:introduction}

\begin{guidance}
The introduction is a high-level overview of the whole work. Explain the motivation and
why the problem matters, how it has been approached so far (briefly; the details belong to
\cref{ch:background}), what is missing and what this thesis contributes, without going into
technical detail. Finish with the structure of the document, one or two sentences per
chapter. Write it, or at least revise it, at the end, when you know exactly what the thesis
contains.
\end{guidance}

\section{Motivation}\label{sec:motivation}

[Describe the context of the problem and why it is relevant: scientific, social, clinical,
economic or industrial impact. Support your claims with references.]

\section{Problem statement}\label{sec:problem}

[State clearly and concisely the problem addressed in this thesis and the gap in the
current knowledge or practice that it aims to fill.]

\section{Contributions}\label{sec:contributions}

[List the main contributions of the thesis. Be concrete: a new method, a dataset, an
empirical comparison, a software tool, a validated prototype\dots]

\section{Structure of the document}\label{sec:structure}

The rest of this document is organised as follows. \Cref{ch:background} reviews the
theoretical background and the state of the art. \Cref{ch:objectives} defines the
objectives of the thesis. \Cref{ch:methods} describes the materials and methods, and
\cref{ch:results} presents the results, which are discussed in \cref{ch:discussion}.
Finally, \cref{ch:conclusions} summarises the conclusions and outlines future work.
